\documentclass[11pt]{article}
\usepackage[margin=1in]{geometry}
\usepackage{booktabs}
\usepackage{enumitem}
\usepackage{hyperref}
\usepackage{xcolor}
\usepackage{array}
\usepackage{longtable}

\definecolor{done}{RGB}{0,128,0}
\definecolor{fix}{RGB}{180,0,0}
\definecolor{check}{RGB}{180,100,0}

\newcommand{\done}{\textcolor{done}{$\checkmark$}}
\newcommand{\fix}{\textcolor{fix}{$\square$}}
\newcommand{\check}{\textcolor{check}{$\square$}}

\title{\textbf{Revision Checklist}\\
\large The Rise of Reddit: How Social Media Affects Belief Formation and Price Discovery\\[0.5em]
\normalsize Management Science Resubmission --- August 2026}
\author{Danqi Hu, Charles M.\ Jones, Siyang Li, Valerie Zhang, Xiaoyan Zhang}
\date{}

\begin{document}
\maketitle
\tableofcontents
\newpage

%----------------------------------------------------------------------
\section{To-Do List}
%----------------------------------------------------------------------

\subsection{Must Fix --- Errors}

\begin{enumerate}[leftmargin=*, label=\fix]

  \item \textbf{Footnote [8] is blank.} The footnote marker appears in
    Section~3.1 but no text follows it. Restore the missing footnote
    content (likely a note about the influence measure or network
    definition).

  \item \textbf{``P.\ Gal'' $\to$ ``P.\ Gao'' in references.} The
    Da, Engelberg, Gao (2011) entry currently reads ``Da, Z.,
    J.~Engelberg, and P.~Gal.\ 2011.'' The author is Pengjie Gao.
    Fix in the reference list.

  \item \textbf{Atmaz (2024) missing from reference list.} Cited in
    Section~4.4 (``Atmaz (2024) shows that persistent demand from
    retail or noise traders can raise stock prices\ldots'') but absent
    from the references. Add the full entry.

  \item \textbf{Engelberg, Reed, and Ringgenberg (2018) missing from
    reference list.} Cited in Section~4.4. The 2012 paper by the same
    authors is present, but 2018 is a different paper
    (``Short-Selling Risk,'' \textit{Journal of Finance}) and is not
    listed. Add the full entry.

  \item \textbf{Formatting artifact in Section~6 header.} A string of
    garbled characters appears just before ``6.1 Technological
    Progress'' in the Word document. Remove.

  \item \textbf{Text--table inconsistency in Section~4.1.} The text
    states that for the low-influence subsample, ``the coefficient on
    RationalTone (t-1) [predicting Na\"{i}veTone] is $-0.004$
    (t-stat$=-0.08$).'' Table~2 Panel~C shows $-0.0004$ [$-0.08$].
    The t-statistic matches but the coefficient in the text is off by
    one decimal place. Correct the text to $-0.0004$.

  \item \textbf{Table~9 note typo: ``other uses'' $\to$ ``other
    users.''} The table note reads ``the separation of celebrity and
    other uses.''

  \item \textbf{``Commentors'' $\to$ ``commenters'' throughout tables.}
    Misspelled as ``commentors'' in the Table~6 Panel~A and Panel~B
    headers, and in the Table~8 note.

\end{enumerate}

\subsection{Potential Inconsistencies --- Review with Coauthors}

\begin{enumerate}[leftmargin=*, label=\check]

  \item \textbf{Table~11: rational tone still predicts retail flow
    during platform-emergence periods.} Section~6.2 states that
    ``predictive relations of Reddit tone are substantially weaker
    during periods when competing platforms are more prominent'' and
    that ``this attenuation also extends to subsequent \ldots trading
    dynamics.'' However, in Table~11 Panel~A (emergence periods), the
    coefficient on RationalTone(t-1) for RetailFlow(t) is $0.0115^{**}$
    ($t=2.16$, Granger $p=3.1\%$) --- still significant at the 5\%
    level. The claim of broad attenuation in retail trading may need
    to be qualified.

  \item \textbf{Table~8 Panel~B: return predictability significant
    even in the low-residual subsample.} The text claims return
    predictability is ``substantially stronger when either the
    user fixed-effect component or the residual component is high.''
    But in Table~8 Panel~B column~IV (low residual), RationalTone
    ($0.0010^{**}$), FanaticTone ($0.0014^{**}$), and Na\"{i}veTone
    ($0.0010^{***}$) are all still significant. Consider softening the
    claim or adding a sentence acknowledging that some predictability
    persists in the low-residual group.

  \item \textbf{Table~9: celebrity short-flow effects are all
    insignificant.} The text focuses on celebrity effects on returns
    and retail flow but does not discuss the null short-flow result.
    All celebrity and other-user tone coefficients for ShortFlow(t)
    are insignificant (Granger $p$-values: 68.6\%, 63.6\%, 57.8\%
    for celebrity; 98.7\%, 74.3\%, 5.4\% for other). Worth at least
    a sentence of acknowledgment.

\end{enumerate}

\subsection{Verified --- No Action Needed}

\begin{enumerate}[leftmargin=*, label=\done]

  \item \textbf{Appendix A present and complete.} LLM classification
    prompt is there, including the full Pedersen (2022) model
    description, the Keith Gill fanatic example, and the rational
    investor (Encore Wire Corporation) example.

  \item \textbf{Appendix B present and complete.} Earnings surprise
    regression table is present with correct coefficients
    (RationalTone: $0.006^{**}$, $t=2.12$; FanaticTone: $-0.002$,
    $t=-0.80$; Na\"{i}veTone: $-0.006^{**}$, $t=-2.60$).

  \item \textbf{Footnote [9] cross-reference to Appendix B.}
    Appears correctly in Section~3.1 and consistent with footnotes
    in the response letters.

  \item \textbf{All other in-text numbers match their tables.}
    Checked coefficients and t-statistics cited in Sections~4.1
    (Tables~2), 4.2 (Table~3), 4.3 (Table~4), 4.4 (Table~5), 5.1
    (Table~6), 5.2 (Table~7), 6.1 (Table~10). All match exactly
    except the Table~2/Section~4.1 discrepancy noted above.

  \item \textbf{Observation counts internally consistent.}
    Main sample: 710,378 (2020--2025). Table~10 pre- and
    post-ChatGPT split: $354{,}975 + 355{,}403 = 710{,}378$. Table~11
    restricted to 2020--2023: $49{,}614 + 436{,}179 = 485{,}793$.
    All consistent.

  \item \textbf{All other cited references present.} Verified: Allen
    et al.\ (2025), Lopez-Lira and Tang (2026), Alsan and Neberai
    (2026), Fedyk et al.\ (2026), Bick et al.\ (2026), Noy and Zhang
    (2023), Cookson et al.\ (2024), Bradshaw et al.\ (2026),
    Hirshleifer et al.\ (2026), Emett et al.\ (2024) are all in the
    reference list.

\end{enumerate}

\newpage
%----------------------------------------------------------------------
\section{Editor Comments}
%----------------------------------------------------------------------

\begin{longtable}{p{0.05\textwidth} p{0.5\textwidth} p{0.38\textwidth}}
\toprule
& \textbf{Comment} & \textbf{Status} \\
\midrule
\endfirsthead
\toprule
& \textbf{Comment} & \textbf{Status} \\
\midrule
\endhead
\bottomrule
\endfoot

\done & Reposition paper as using Pedersen (2022) as \textit{conceptual motivation} rather than a direct test. &
  Addressed. Abstract, Section~1, Section~2 (renamed ``Conceptual Framework and Empirical Implications''), and Section~7 revised. \\[6pt]

\done & Differentiate types of social media interactions (anonymous vs.\ celebrity posters). &
  Addressed. Section~5.4, Table~9. Celebrity users defined as those with average PageRank above weekly 90th percentile. \\[6pt]

\done & Relate to technological progress. &
  Addressed. Section~6.1, Table~10. ChatGPT release (Nov.\ 30, 2022) used as cutoff; return predictability declines post-ChatGPT. \\[6pt]

\done & Emergence of competing platforms. &
  Addressed. Section~6.2, Table~11. Relative activity on StockTwits, X, and Seeking Alpha used to identify high-competition periods. \\

\end{longtable}

%----------------------------------------------------------------------
\section{Reviewer 1 Comments}
%----------------------------------------------------------------------

\begin{longtable}{p{0.05\textwidth} p{0.5\textwidth} p{0.38\textwidth}}
\toprule
& \textbf{Comment} & \textbf{Status} \\
\midrule
\endfirsthead
\toprule
& \textbf{Comment} & \textbf{Status} \\
\midrule
\endhead
\bottomrule
\endfoot

\done & \textbf{Comment 1.} PVAR insufficient to distinguish belief spillover from autocorrelation of individual investors (dynamic reshuffling concern). &
  Addressed. New LLM-based classification (GPT-5 on full posting histories). Section~4.1 reframed to clarify group-level interpretation. \\[6pt]

\done & \textbf{Comment 2.} Investor classification ignores alignment with fundamentals. &
  Addressed. LLM uses Pedersen (2022) definitions directly. Earnings surprise validation added: only rational tone positively predicts ($t=2.12$). Section~3.1. \\[6pt]

\done & \textbf{Comment 3.} Rational tone no better than na\"{i}ve at predicting returns. &
  Addressed. Updated results: RationalTone $t=2.53$, FanaticTone $t=2.80$, Na\"{i}veTone $t=3.94$ --- all significant, similar magnitude. Table~3. \\[6pt]

\done & \textbf{Comment 4.} ``Influence'' variable may proxy for attention, not genuine social influence. &
  Addressed. PageRank-based centrality measure introduced (Section~5.2). Correlation between centrality and attention is only 0.25. \\[6pt]

\done & \textbf{Comment 5.} Rational/fanatic tone does not predict retail trading flows. &
  Addressed. Updated results: RationalTone $t=6.80$, FanaticTone $t=6.24$ --- both significant. Table~4. \\[6pt]

\done & \textbf{Other 1.} Typo: ``do not predict'' should be ``insignificant.'' &
  Addressed. Corrected in Section~4.4. \\[6pt]

\done & \textbf{Other 2.} Short sellers ``detecting influence'' claim is speculative. &
  Addressed. Statement toned down in abstract. \\[6pt]

\done & \textbf{Other 3.} Reddit users and retail investors likely overlap. &
  Addressed. Explicitly acknowledged in Section~4.3. \\

\end{longtable}

%----------------------------------------------------------------------
\section{Reviewer 2 Comments}
%----------------------------------------------------------------------

\begin{longtable}{p{0.05\textwidth} p{0.5\textwidth} p{0.38\textwidth}}
\toprule
& \textbf{Comment} & \textbf{Status} \\
\midrule
\endfirsthead
\toprule
& \textbf{Comment} & \textbf{Status} \\
\midrule
\endhead
\bottomrule
\endfoot

\done & \textbf{C1.} Empirical proxies do not match theoretical rational/fanatic types; na\"{i}ve tone is the strongest return predictor; rational tone fails to predict returns. &
  Addressed. LLM reclassification; rational tone now predicts returns at $t=2.53$; earnings surprise test; long-horizon portfolio analysis (52 weeks). Sections~3.1, 4.2, 4.5. \\[6pt]

\done & \textbf{C2.} Influence measure is endogenous (stock $\times$ time; reverse causality concern). &
  Addressed. Three alternative measures: (1) lagged cross-stock average commenters (Section~5.1); (2) pre-period PageRank centrality (Section~5.2); (3) fixed-effect decomposition into persistent and transient components (Section~5.3). \\[6pt]

\done & \textbf{C3.} Retail flow is not a valid proxy for na\"{i}ve trading; tone$\to$flow$\to$price chain not supported. &
  Addressed. Explicitly acknowledged in Section~4.3. Bidirectional Granger causality reported (tone $\to$ flow and flow $\to$ tone both significant). \\[6pt]

\done & \textbf{C4.} No momentum-reversal pattern, contradicting Pedersen (2022). &
  Addressed. Section~4.5, Figure~6. Na\"{i}ve tone reverses to $-1.6\%$ by week~52; rational tone accumulates to $+4\%$; fanatic tone dissipates near zero. \\

\end{longtable}

\end{document}
